\chapter{Limites e continuidade}\label{ch:g12:limcont}

Este capítulo estende a noção de limite das \omterm{def:g12:seq:sequence}{sequências} para as funções de
uma variável real e apresenta a \omterm{def:g12:limcont:continuity}{continuidade}. O resultado central é o
teorema do valor intermediário, que transforma informação qualitativa
(``$f$ é \omterm{def:g12:limcont:continuity}{contínua} e muda de sinal'') na existência de soluções de
\omterm{def:g10:algebra:equation}{equações}.

\section{Limites de funções}

Ao longo do capítulo, $f$ é uma \omterm{def:g11:func:function}{função} definida em um \omterm{def:g10:numbers:interval}{intervalo} ou em uma
\omterm{def:g10:numbers:interunion}{união} de \omterm{def:g10:numbers:interval}{intervalos} $D \subseteq \R$.

\begin{definition}[Limite no infinito]\label{def:g12:limcont:liminf}
Seja $f$ definida em um \omterm{def:g10:numbers:interval}{intervalo} $\intco{a}{+\infty}$ e seja
$\ell \in \R$. Dizemos que $f$ \emph{tende a $\ell$ em
$+\infty$}\index{limite!de uma função}, o que se escreve
$\lim\limits_{x\to+\infty} f(x) = \ell$, se todo \omterm{def:g10:numbers:interval}{intervalo} aberto que
contém $\ell$ contém todos os valores $f(x)$ para $x$ suficientemente
grande: para todo $\varepsilon > 0$ existe $A$ tal que, para todo
$x \geq A$, $\abs{f(x) - \ell} \leq \varepsilon$.

Dizemos que $f$ tende a $+\infty$ em $+\infty$ se, para todo $M \in \R$,
existe $A$ tal que $f(x) \geq M$ para todo $x \geq A$. Os limites em
$-\infty$ e os limites iguais a $-\infty$ definem-se de modo análogo.
\end{definition}

\begin{definition}[Limite em um ponto]\label{def:g12:limcont:limpoint}
Seja $a \in \R$ e seja $f$ definida perto de $a$ (exceto, talvez, em
$a$). Dizemos que $f$ tende a $\ell \in \R$ em $a$, o que se escreve
$\lim\limits_{x \to a} f(x) = \ell$, se, para todo $\varepsilon > 0$,
existe $\delta > 0$ tal que, para todo $x \in D$ com
$\abs{x - a} \leq \delta$, $\abs{f(x) - \ell} \leq \varepsilon$.

Dizemos que $f$ tende a $+\infty$ em $a$ se, para todo $M$, existe
$\delta > 0$ tal que $f(x) \geq M$ sempre que $x \in D$ e
$\abs{x - a} \leq \delta$. Os limites laterais ($x \to a^+$,
$x \to a^-$) restringem $x$ a $x > a$ ou $x < a$.
\end{definition}

\begin{example}
$\lim\limits_{x\to+\infty} \frac{1}{x} = 0$,
$\lim\limits_{x\to 0^+} \frac{1}{x} = +\infty$,
$\lim\limits_{x\to 0^-} \frac{1}{x} = -\infty$. A \omterm{def:g11:func:function}{função}
$x \mapsto \frac1x$ não tem limite (bilateral) em $0$.
\end{example}

\begin{definition}[Assíntotas]\label{def:g12:limcont:asymptote}
A reta $y = \ell$ é uma \emph{assíntota horizontal}\index{assíntota} da
curva de $f$ em $+\infty$ (resp.\ $-\infty$) se
$\lim_{x\to+\infty} f(x) = \ell$ (resp.\ em $-\infty$). A reta $x = a$ é
uma \emph{assíntota vertical}\index{assíntota} se $f$ tem um limite
lateral infinito em $a$.
\end{definition}

\begin{omfigure}
\begin{tikzpicture}
\begin{axis}[omaxis,
  width=8.8cm, height=6cm,
  xmin=-2.9, xmax=4.9, ymin=-5.2, ymax=9.2,
  xtick={1}, ytick={2}]
  \draw[omProp, dashed] (axis cs:-2.9,2) -- (axis cs:4.9,2);
  \draw[omProp, dashed] (axis cs:1,-5.2) -- (axis cs:1,9.2);
  \addplot[omDef, thick, domain=-2.8:0.855, samples=60] {2 + 1/(x-1)};
  \addplot[omDef, thick, domain=1.145:4.8, samples=60] {2 + 1/(x-1)};
\end{axis}
\end{tikzpicture}

{\small A curva de $f(x) = 2 + \frac{1}{x-1}$ com sua \omterm{def:g12:limcont:asymptote}{assíntota
horizontal} $y = 2$ (em $\pm\infty$) e sua \omterm{def:g12:limcont:asymptote}{assíntota vertical} $x = 1$.}
\end{omfigure}

\begin{proposition}[Operações com limites]\label{prop:g12:limcont:operations}
As regras da \cref{prop:g12:seq:operations} (somas, produtos, quocientes)
valem literalmente para limites de funções, em um ponto ou no infinito,
com as mesmas quatro formas indeterminadas $(+\infty)+(-\infty)$,
$0\times\infty$, $\frac{\infty}{\infty}$, $\frac{0}{0}$.
\end{proposition}

\begin{proof}
As demonstrações são idênticas às do caso das \omterm{def:g12:seq:sequence}{sequências}, trocando ``para
$n \geq N$'' por ``para $x \geq A$'' ou ``para
$\abs{x-a} \leq \delta$''.
\end{proof}

\begin{theorem}[Limite de uma composta]\label{thm:g12:limcont:composition}
Sejam $f, g$ funções e $a, b, c$ \omterm{def:g10:numbers:sets}{números reais} ou $\pm\infty$. Se
$\lim\limits_{x \to a} f(x) = b$ e $\lim\limits_{y \to b} g(y) = c$,
então
\[
\lim_{x \to a} g\bigl(f(x)\bigr) = c .
\]
O mesmo enunciado vale com uma \omterm{def:g12:seq:sequence}{sequência} no lugar de $f$: se
$u_n \to b$ e $\lim_{y\to b} g(y) = c$, então $g(u_n) \to c$.
\end{theorem}

\begin{proof}
Tome $b, c$ finitos (os demais casos são análogos). Seja
$\varepsilon > 0$. Existe $\delta > 0$ tal que $\abs{y - b} \leq \delta$
implica $\abs{g(y) - c} \leq \varepsilon$; e existe $\delta' > 0$ tal que
$\abs{x - a} \leq \delta'$ implica $\abs{f(x) - b} \leq \delta$. Encadear
os dois dá $\abs{g(f(x)) - c} \leq \varepsilon$ para
$\abs{x - a} \leq \delta'$.
\end{proof}

\begin{theorem}[Comparação e confronto]\label{thm:g12:limcont:squeeze}
Sejam $f, g, h$ funções e $a \in \R \cup \{\pm\infty\}$.
\begin{enumerate}
  \item Se $f \leq g$ perto de $a$ e $f(x) \to +\infty$ quando
        $x \to a$, então $g(x) \to +\infty$.
  \item Se $f \leq g \leq h$ perto de $a$ e $f, h$ tendem ao mesmo limite
        finito $\ell$ em $a$, então $g(x) \to \ell$ quando $x \to a$.
\end{enumerate}
\end{theorem}

\begin{proof}
Mesmo argumento do caso das \omterm{def:g12:seq:sequence}{sequências} (\cref{thm:g12:seq:squeeze}).
\end{proof}

\begin{method}[Calcular limites na prática]\label{met:g12:limcont:limits}
\begin{itemize}
  \item Polinômios e funções racionais em $\pm\infty$: coloque em
        evidência a maior potência de $x$; o limite é o da razão dos
        termos dominantes.
  \item $\frac{0}{0}$ em um ponto $a$: fatore $(x-a)$ no numerador e no
        denominador, ou reconheça uma \omterm{def:g11:deriv:derivative}{derivada}
        $\lim_{x\to a}\frac{f(x)-f(a)}{x-a} = f'(a)$
        (\cref{ch:g12:deriv}).
  \item Raízes quadradas: multiplique pela expressão conjugada.
  \item Fatores oscilantes ($\cos$, $\sin$): use o teorema do confronto.
\end{itemize}
\end{method}

\section{Continuidade}

\begin{definition}[Continuidade]\label{def:g12:limcont:continuity}
Seja $f$ definida em um \omterm{def:g10:numbers:interval}{intervalo} $I$ e seja $a \in I$. A \omterm{def:g11:func:function}{função} $f$ é
\emph{contínua em $a$}\index{continuidade} se
$\lim\limits_{x \to a} f(x) = f(a)$. Ela é \emph{contínua em
$I$}\index{contínua} se for contínua em todo ponto de $I$.
\end{definition}

Intuitivamente, o \omterm{def:g10:functions:graph}{gráfico} de uma \omterm{def:g11:func:function}{função} \omterm{def:g12:limcont:continuity}{contínua} em um \omterm{def:g10:numbers:interval}{intervalo} pode ser
traçado sem levantar o lápis.

\begin{proposition}[Continuidade das funções usuais]\label{prop:g12:limcont:usual}
Polinômios, funções racionais, $x \mapsto \sqrt{x}$, $x \mapsto \abs{x}$,
$\cos$, $\sin$, $\exp$ e $\ln$ são \omterm{def:g12:limcont:continuity}{contínuas} em todo \omterm{def:g10:numbers:interval}{intervalo} contido em
seu \omterm{def:g11:func:function}{domínio}. Somas, produtos, quocientes (onde definidos) e composições de
funções \omterm{def:g12:limcont:continuity}{contínuas} são \omterm{def:g12:limcont:continuity}{contínuos}.
\end{proposition}

\begin{proof}[Demonstração parcial]
A estabilidade por somas, produtos, quocientes e composição segue dos
teoremas correspondentes sobre limites. A \omterm{def:g12:limcont:continuity}{continuidade} de $x \mapsto x$ e
das funções constantes é imediata pela definição, de modo que polinômios e
funções racionais são \omterm{def:g12:limcont:continuity}{contínuos}. A \omterm{def:g12:limcont:continuity}{continuidade} das demais funções usuais
é demonstrada nos capítulos em que elas são construídas, ou admitida.
\end{proof}

\begin{example}
A \emph{\omterm{def:g11:func:function}{função} piso} $x \mapsto \floor{x}$ (o maior inteiro $\leq x$) não
é \omterm{def:g12:limcont:continuity}{contínua} em nenhum inteiro $n$: seu limite à esquerda ali é $n - 1$,
enquanto $\floor{n} = n$.
\end{example}

\begin{omfigure}
\begin{tikzpicture}
\begin{axis}[omaxis,
  width=7.2cm, height=5.6cm,
  xmin=-1.9, xmax=3.5, ymin=-2.5, ymax=2.7,
  xtick={-1,1,2,3}, ytick={-2,-1,1,2}]
  \draw[omDef, thick] (axis cs:-1.5,-2) -- (axis cs:-1,-2);
  \draw[omDef, thick] (axis cs:-1,-1) -- (axis cs:0,-1);
  \draw[omDef, thick] (axis cs:0,0) -- (axis cs:1,0);
  \draw[omDef, thick] (axis cs:1,1) -- (axis cs:2,1);
  \draw[omDef, thick] (axis cs:2,2) -- (axis cs:3,2);
  \fill[omDef] (axis cs:-1,-1) circle (1.5pt);
  \fill[omDef] (axis cs:0,0) circle (1.5pt);
  \fill[omDef] (axis cs:1,1) circle (1.5pt);
  \fill[omDef] (axis cs:2,2) circle (1.5pt);
  \draw[omDef, fill=white] (axis cs:-1,-2) circle (1.5pt);
  \draw[omDef, fill=white] (axis cs:0,-1) circle (1.5pt);
  \draw[omDef, fill=white] (axis cs:1,0) circle (1.5pt);
  \draw[omDef, fill=white] (axis cs:2,1) circle (1.5pt);
  \draw[omDef, fill=white] (axis cs:3,2) circle (1.5pt);
\end{axis}
\end{tikzpicture}

{\small A \omterm{def:g11:func:function}{função} piso salta em cada inteiro: o valor (bolinha cheia)
difere do limite à esquerda (bolinha vazia).}
\end{omfigure}

\begin{proposition}[Sequências e continuidade]\label{prop:g12:limcont:seqcont}
Se $u_n \to \ell$ e $f$ é \omterm{def:g12:limcont:continuity}{contínua} em $\ell$, então
$f(u_n) \to f(\ell)$. Em particular, se uma \omterm{def:g12:seq:sequence}{sequência} definida por
$u_{n+1} = f(u_n)$ \omterm{def:g12:seq:limit}{converge} para $\ell$ e $f$ é \omterm{def:g12:limcont:continuity}{contínua} em $\ell$, então
$f(\ell) = \ell$.
\end{proposition}

\begin{proof}
A primeira afirmação é o \cref{thm:g12:limcont:composition} aplicado à
\omterm{def:g12:seq:sequence}{sequência} $(u_n)$. Para a segunda, passe ao limite nos dois lados de
$u_{n+1} = f(u_n)$: o lado esquerdo tende a $\ell$, o direito a
$f(\ell)$, e os limites são únicos.
\end{proof}

\section{O teorema do valor intermediário}

\begin{theorem}[Teorema do valor intermediário]\label{thm:g12:limcont:ivt}
\index{teorema do valor intermediário}
Seja $f$ \omterm{def:g12:limcont:continuity}{contínua} em $\intcc{a}{b}$ e seja $k$ um número qualquer entre
$f(a)$ e $f(b)$. Então existe $c \in \intcc{a}{b}$ tal que $f(c) = k$.
\end{theorem}

\begin{omfigure}
\begin{tikzpicture}
\begin{axis}[omaxis,
  width=8.8cm, height=6cm,
  xmin=-0.4, xmax=4.4, ymin=-0.5, ymax=4.7,
  xtick={0.5,3.19,3.8}, xticklabels={$a$,$c$,$b$},
  ytick={1.26,2.5,3.98}, yticklabels={$f(a)$,$k$,$f(b)$}]
  \addplot[omProp, dashed, domain=0:3.19] {2.5};
  \addplot[omDef, thick, domain=0.5:3.8] {0.08*x^3 - 0.5*x + 1.5};
  \draw[black, dashed, thin] (axis cs:0.5,0) -- (axis cs:0.5,1.26)
    -- (axis cs:0,1.26);
  \draw[black, dashed, thin] (axis cs:3.8,0) -- (axis cs:3.8,3.98)
    -- (axis cs:0,3.98);
  \draw[black, dashed, thin] (axis cs:3.19,2.5) -- (axis cs:3.19,0);
  \fill (axis cs:0.5,1.26) circle (1.5pt);
  \fill (axis cs:3.8,3.98) circle (1.5pt);
  \fill (axis cs:3.19,2.5) circle (1.5pt);
\end{axis}
\end{tikzpicture}

{\small O teorema do valor intermediário: uma curva \omterm{def:g12:limcont:continuity}{contínua} que liga
$(a, f(a))$ a $(b, f(b))$ tem de cruzar toda reta horizontal $y = k$
situada entre $f(a)$ e $f(b)$.}
\end{omfigure}

\begin{proof}
Substituindo $f$ por $f - k$ (e por $k - f$ se preciso), podemos supor
$f(a) \leq 0 \leq f(b)$ e procurar um zero de $f$. Construímos duas
\omterm{def:g12:seq:sequence}{sequências} por \emph{dicotomia}\index{dicotomia}: ponha $a_0 = a$,
$b_0 = b$; dados $a_n \leq b_n$ com $f(a_n) \leq 0 \leq f(b_n)$, seja
$m = \frac{a_n + b_n}{2}$ e defina
\[
(a_{n+1}, b_{n+1}) =
\begin{cases}
(a_n, m) & \text{se } f(m) \geq 0,\\
(m, b_n) & \text{se } f(m) < 0.
\end{cases}
\]
Nos dois casos $f(a_{n+1}) \leq 0 \leq f(b_{n+1})$, a \omterm{def:g12:seq:sequence}{sequência} $(a_n)$ é
\omterm{def:g12:seq:monotonic}{crescente}, $(b_n)$ é \omterm{def:g10:functions:variations}{decrescente} e $b_n - a_n = \frac{b-a}{2^n} \to 0$:
as \omterm{def:g12:seq:sequence}{sequências} são adjacentes (veja o \cref{exo:g12:seq:10}) e convergem
para um limite comum $c \in \intcc{a}{b}$. Pela \omterm{def:g12:limcont:continuity}{continuidade},
$f(a_n) \to f(c)$ e $f(b_n) \to f(c)$; passar ao limite em
$f(a_n) \leq 0$ e $f(b_n) \geq 0$ dá $f(c) \leq 0 \leq f(c)$, logo
$f(c) = 0$.
\end{proof}

\begin{remark}
A demonstração por \omterm{thm:g12:limcont:ivt}{dicotomia} é um algoritmo: ela fornece \omterm{def:g10:numbers:approx}{aproximações}
arbitrariamente precisas de uma solução, dividindo o erro pela metade a
cada passo. É assim que um computador resolve $f(x) = 0$ sabendo apenas
avaliar $f$.
\end{remark}

\begin{theorem}[Teorema da bijeção]\label{thm:g12:limcont:bijection}
Seja $f$ \omterm{def:g12:limcont:continuity}{contínua} e \emph{estritamente \omterm{def:g12:seq:monotonic}{monótona}} em um \omterm{def:g10:numbers:interval}{intervalo} $I$.
Então, para todo $k$ estritamente entre os limites de $f$ nos extremos de
$I$, a \omterm{def:g10:algebra:equation}{equação} $f(x) = k$ tem uma \emph{única} solução em $I$.
\end{theorem}

\begin{proof}
Existência: escolha $a, b \in I$ com $k$ entre $f(a)$ e $f(b)$ (possível,
já que os valores de $f$ se aproximam dos seus limites nos extremos) e
aplique o teorema do valor intermediário em $\intcc{a}{b}$. Unicidade: se
$x < y$, a monotonicidade estrita dá $f(x) \neq f(y)$, de modo que dois
pontos distintos não podem ser ambos soluções.
\end{proof}

\begin{method}[Resolver $f(x) = k$ com o teorema da bijeção]\label{met:g12:limcont:bijection}
Para demonstrar que $f(x) = k$ tem exatamente uma solução em um
\omterm{def:g10:numbers:interval}{intervalo}:
\begin{enumerate}
  \item estabeleça um \omterm{def:g10:functions:table}{quadro de variação} de $f$ (sinal de $f'$, veja o
        \cref{ch:g12:deriv});
  \item em cada \omterm{def:g10:numbers:interval}{intervalo} em que $f$ é \omterm{def:g12:limcont:continuity}{contínua} e estritamente \omterm{def:g12:seq:monotonic}{monótona},
        compare $k$ com os valores ou limites nos extremos;
  \item aplique o teorema da bijeção em cada um desses \omterm{def:g10:numbers:interval}{intervalos} e conte
        as soluções.
\end{enumerate}
Uma calculadora ou a \omterm{thm:g12:limcont:ivt}{dicotomia} localizam depois cada solução com a
precisão que se quiser.
\end{method}

\begin{example}\label{ex:g12:limcont:cubic}
A \omterm{def:g10:algebra:equation}{equação} $x^3 + x = 1$ tem uma única solução real. De fato,
$f(x) = x^3 + x$ é \omterm{def:g12:limcont:continuity}{contínua} e estritamente \omterm{def:g12:seq:monotonic}{crescente} em $\R$ (soma de
funções estritamente \omterm{def:g12:seq:monotonic}{crescentes}), $\lim_{-\infty} f = -\infty$ e
$\lim_{+\infty} f = +\infty$, de modo que o teorema da bijeção se aplica
com $k = 1$. Como $f(0.6) = 0.816 < 1$ e $f(0.7) = 1.043 > 1$, a solução
está em $\intoo{0.6}{0.7}$.
\end{example}

\section{Exercícios}

\begin{exercise}[$\star$]\label{exo:g12:limcont:1}
Calcule os limites seguintes:
\[
\lim_{x\to+\infty} \frac{2x^3 - x + 1}{5x^3 + x^2}, \qquad
\lim_{x\to-\infty} \bigl(x^2 + 3x - 1\bigr), \qquad
\lim_{x\to 2^+} \frac{x+1}{x-2}, \qquad
\lim_{x\to+\infty} \frac{\cos x}{x}.
\]
\end{exercise}

\begin{exercise}[$\star$]\label{exo:g12:limcont:2}
Seja $f(x) = \dfrac{2x^2 - 3x + 1}{x - 1}$, definida para $x \neq 1$.
\begin{enumerate}
  \item Mostre que $f(x) = 2x - 1$ para todo $x \neq 1$.
  \item Deduza $\lim\limits_{x \to 1} f(x)$. É possível estender $f$ a uma
        \omterm{def:g11:func:function}{função} \omterm{def:g12:limcont:continuity}{contínua} em $1$?
\end{enumerate}
\end{exercise}

\begin{exercise}[$\star$]\label{exo:g12:limcont:3}
Determine as \omterm{def:g12:limcont:asymptote}{assíntotas} da curva de $f(x) = \dfrac{3x - 1}{x + 2}$ e
esboce a curva.
\end{exercise}

\begin{exercise}[$\star\star$]\label{exo:g12:limcont:4}
Calcule
\[
\lim_{x\to+\infty} \bigl(\sqrt{x^2 + x} - x\bigr)
\qquad\text{e}\qquad
\lim_{x\to 0} \frac{\sqrt{1+x} - 1}{x}.
\]
\end{exercise}

\begin{exercise}[$\star\star$]\label{exo:g12:limcont:5}
Seja $f(x) = x + \sin x$.
\begin{enumerate}
  \item Mostre que $f(x) \to +\infty$ quando $x \to +\infty$, embora $f$
        não seja \omterm{def:g12:seq:monotonic}{monótona} em nenhum \omterm{def:g10:numbers:interval}{intervalo} $\intco{A}{+\infty}$.
  \item Mostre que $g(x) = \dfrac{x + \sin x}{x - \sin x}$ tende a $1$ em
        $+\infty$.
\end{enumerate}
\end{exercise}

\begin{exercise}[$\star\star$]\label{exo:g12:limcont:6}
Mostre que a \omterm{def:g10:algebra:equation}{equação} $x^3 - 3x + 1 = 0$ tem exatamente três soluções
reais e dê, para cada uma, um \omterm{def:g10:numbers:interval}{intervalo} de comprimento $1$ que a contenha.
(Sugestão: estude a variação de $f(x) = x^3 - 3x + 1$; sua \omterm{def:g11:deriv:derivative}{derivada} é
$3x^2 - 3$.)
\end{exercise}

\begin{exercise}[$\star\star$]\label{exo:g12:limcont:7}
Seja $f \colon \intcc{0}{1} \to \intcc{0}{1}$ \omterm{def:g12:limcont:continuity}{contínua}. Mostre que $f$ tem
um \emph{\omterm{pb:g10:functions:1}{ponto fixo}}: existe $c \in \intcc{0}{1}$ com $f(c) = c$.
(Sugestão: aplique o teorema do valor intermediário a $g(x) = f(x) - x$.)
\end{exercise}

\begin{exercise}[$\star\star\star$]\label{exo:g12:limcont:8}
Uma caminhante sobe uma trilha na montanha, partindo às 8h e chegando às
20h. No dia seguinte ela desce a mesma trilha, de novo partindo às 8h e
chegando às 20h. Mostre que existe um horário do dia em que ela está
exatamente no mesmo ponto da trilha nos dois dias. (Sugestão: introduza a
diferença das duas funções de posição.)
\end{exercise}

\begin{exercise}[$\star\star\star$]\label{exo:g12:limcont:9}
Seja $f(x) = x^5 + x^3 + x - 1$.
\begin{enumerate}
  \item Mostre que $f$ tem um único zero real $\alpha$ e que
        $\alpha \in \intoo{0}{1}$.
  \item Descreva um procedimento de \omterm{thm:g12:limcont:ivt}{dicotomia} que calcule $\alpha$ com
        precisão $10^{-6}$ e dê o número de iterações necessárias partindo
        de $\intcc{0}{1}$.
\end{enumerate}
\end{exercise}

\section{Problema: Não se atravessa um rio sem se molhar}

\begin{problem}[{Problema de fim de semana --- o teorema do valor
intermediário força pontos fixos, antípodas de mesma temperatura e
mesas firmes: três teoremas de aparência impossível a partir de uma
única ideia}]
\label{pb:g12:limcont:1}
Amasse um mapa do Brasil e deixe-o cair em qualquer ponto do
território brasileiro: algum ponto do mapa fica \emph{exatamente}
acima do lugar que representa. A cada instante, dois pontos
diametralmente opostos do equador têm \emph{exatamente} a mesma
temperatura. E uma mesa quadrada bamba, em chão irregular, sempre
pode ser firmada \emph{girando-a}. Três truques de festa, um
teorema: uma grandeza \omterm{def:g12:limcont:continuity}{contínua} não passa de negativa a positiva
sem se anular (\cref{thm:g12:limcont:ivt}). Este problema treina
limites, constrói em seguida as três maravilhas --- e sabe
exatamente o que o teorema não promete.

\medskip
\noindent\textbf{Parte I --- Fluência com limites.}
\begin{enumerate}
  \item Calcule $\lim_{x \to +\infty} \dfrac{3x^2 - x}{x^2 + 5}$,
        \ $\lim_{x \to +\infty} \dfrac{2x + 1}{x^2 + 1}$ \ e
        $\lim_{x \to +\infty} (x^3 - x^2)$
        (\cref{met:g12:limcont:limits}).
  \item Dê as \omterm{def:g12:limcont:asymptote}{assíntotas} de $f(x) = \dfrac{3x - 1}{x + 2}$
        (\cref{def:g12:limcont:asymptote}), calculando os limites
        relevantes.
  \item Calcule $\lim_{x \to 3} \dfrac{x^2 - 9}{x - 3}$ (fatore)
        e $\lim_{x \to 0} \dfrac{\sqrt{x + 1} - 1}{x}$
        (conjugado).
  \item Pelo teorema do confronto
        (\cref{thm:g12:limcont:squeeze}): calcule
        $\lim_{x \to 0} x \sin\frac1x$ e
        $\lim_{x \to +\infty} \frac{\sin x}{x}$.
  \item Por composição e termos dominantes:
        $\lim_{x \to +\infty} \dfrac{\sqrt{4x^2 + x}}{x}$.
\end{enumerate}

\medskip
\noindent\textbf{Parte II --- Raízes, certificadas.}
\begin{enumerate}[resume]
  \item Mostre que a \omterm{def:g10:algebra:equation}{equação} $x^3 + x = 5$ tem pelo menos uma
        solução em $\intcc{1}{2}$.
  \item Mostre que a solução é única em $\R$
        (\cref{thm:g12:limcont:bijection}).
  \item Demonstre o clássico: \emph{todo polinômio de grau ímpar
        tem pelo menos uma \omterm{def:g11:quad:discriminant}{raiz} real}. (Escreva o argumento para
        $p(x) = x^3 + a x^2 + b x + c$: calcule os dois limites
        infinitos pelo termo dominante e aplique em seguida o
        teorema do valor intermediário em um intervalo grande.)
  \item Localizando a \omterm{def:g11:quad:discriminant}{raiz} da questão 6 por \omterm{thm:g12:limcont:ivt}{dicotomia}: quantas
        divisões ao meio de $\intcc{1}{2}$ garantem um erro
        abaixo de $10^{-6}$ (\cref{exo:g12:limcont:9})? Compare
        com a duplicação de dígitos do método de Newton
        (\cref{pb:g11:deriv:1}): de quantos passos ele precisaria,
        aproximadamente?
  \item O contraexemplo que guarda o teorema: $g(x) = \frac1x$
        satisfaz $g(-1) = -1 < 0$ e $g(1) = 1 > 0$ e, no entanto,
        nunca se anula. Que hipótese do teorema falha --- e o que
        o exemplo ensina sobre conferir hipóteses?
\end{enumerate}

\medskip
\noindent\textbf{Parte III --- Três teoremas de aparência
impossível.}
\begin{enumerate}[resume]
  \item O teorema do \omterm{pb:g10:functions:1}{ponto fixo} no segmento: se
        $f : \intcc{0}{1} \to \intcc{0}{1}$ é \omterm{def:g12:limcont:continuity}{contínua},
        demonstre que algum $c$ satisfaz $f(c) = c$. (Estude
        $g(x) = f(x) - x$ nas duas extremidades.)
  \item Traduza a questão 11 para o mapa amassado: o que faz o
        papel de $f$, por que ele cai em $\intcc{0}{1}$
        (figurativamente: dentro do país) e qual é o ponto que
        fica exatamente acima da própria localização? (Uma
        dimensão é admitida em silêncio --- o enunciado honesto
        em duas dimensões é o teorema de Brouwer, dos volumes de
        graduação.)
  \item O equador: seja $T$ uma temperatura \omterm{def:g12:limcont:continuity}{contínua} sobre um
        círculo, vista como uma \omterm{def:g11:func:function}{função} $2\pi$-periódica com
        $T(0) = T(2\pi)$. Ponha $g(t) = T(t) - T(t + \pi)$.
        Calcule $g(0) + g(\pi)$, deduza que $g$ se anula em algum
        ponto de $\intcc{0}{\pi}$ e enuncie o teorema sobre
        pontos antípodas.
  \item A mesa bamba: uma mesa quadrada perfeitamente rígida
        apoia-se em um chão irregular \omterm{def:g12:limcont:continuity}{contínuo} (sem degraus!);
        três pés sempre tocam o chão e um fica no ar a uma altura
        $h(\theta)$ que depende continuamente do ângulo de
        rotação $\theta$ da mesa. Explique por que girar um
        quarto de volta força a folga do pé suspenso a mudar de
        sinal --- e conclua que algum ângulo $\theta^*$ firma os
        quatro pés. (Esse argumento foi publicado como teorema
        legítimo em 2005; os donos de café já conheciam o truque
        antes.)
  \item A caminhante do \cref{exo:g12:limcont:8} usou o mesmo
        esquema. Enuncie em três linhas o modelo comum às
        questões 11, 13, 14 e à caminhante --- construir que
        \omterm{def:g11:func:function}{função}, verificar o quê nas extremidades, invocar o quê.
  \item O que o teorema do valor intermediário \emph{não} dá: a
        existência veio de graça, mas que duas perguntas ele
        deixa em aberto, e quais ferramentas deste problema as
        respondem (questões 7 e 9)?
\end{enumerate}

\medskip
\noindent\textbf{Parte IV --- Retratos com \omterm{def:g12:limcont:asymptote}{assíntotas}.}
\begin{enumerate}[resume]
  \item Reescreva $f(x) = \dfrac{3x - 1}{x + 2}$ como
        $3 - \dfrac{7}{x + 2}$ e use isso para descrever a curva
        por completo: \omterm{def:g12:limcont:asymptote}{assíntotas}, variação de cada lado, esboço.
        (Uma \omterm{def:g11:func:reference}{hipérbole} deslocada --- a curva do custo médio do
        \cref{pb:g10:reffunc:1} também era uma.)
  \item A \omterm{def:g11:func:function}{função} $f(x) = \dfrac{x^2 - 1}{x - 1}$ não está
        definida em $1$. Que valor atribuído a $f(1)$ a torna
        \omterm{def:g12:limcont:continuity}{contínua} ali --- e por que os analistas chamam um ponto
        assim de descontinuidade \emph{removível}?
  \item A \omterm{def:g11:func:function}{função} piso $x \mapsto \lfloor x \rfloor$ (o maior
        inteiro $\leq x$): onde ela é \omterm{def:g12:limcont:continuity}{contínua}, onde salta, e por
        que a conclusão do teorema do valor intermediário falha
        para ela (encontre $a < b$ com
        $\lfloor a \rfloor < \frac12 < \lfloor b \rfloor$ e, no
        entanto, sem solução de
        $\lfloor x \rfloor = \frac12$)?
  \item Final --- as duas faces da \omterm{def:g12:limcont:continuity}{continuidade}: a promessa local
        (valor $=$ limite: nada de teletransporte) e o poder
        global (o teorema do valor intermediário: todo valor
        intermediário é visitado). Classifique as quatro
        maravilhas da Parte III na face certa e dê ao lema da
        travessia do rio seu nome matemático preciso.
\end{enumerate}
\end{problem}
